\documentclass[jou,11pt,donotrepeattitle]{apa7}
% APA 7th edition. Compilar:  latexmk -pdf criterio_snr.tex
\usepackage[american]{babel}
\usepackage{csquotes}
\usepackage[style=apa,backend=biber]{biblatex}
\DeclareLanguageMapping{american}{american-apa}
\addbibresource{referencias.bib}
\usepackage{amsmath,amssymb}
\usepackage{booktabs}
\usepackage{graphicx}
\usepackage{siunitx}

\title{Stopping Criteria Below the Signal-to-Noise Floor: Window Length,
Not Tolerance, Governs Convergence Detection in Architecture Comparisons}
\shorttitle{Stopping Criteria Below the Noise Floor}
\leftheader{Speranza}

\author{Maximiliano Speranza}
\affiliation{Independent Researcher, Argentina}

\authornote{\addORCIDlink{Maximiliano Speranza}{0009-0005-0413-8554}
This work was carried out independently, without institutional affiliation, supervision, or
funding. The author is an undergraduate student of the University Technical Degree in
Programming (\emph{Tecnicatura Universitaria en Programaci\'on}, TUP) at Universidad
Tecnol\'ogica Nacional (UTN), Argentina, and a graduate of its Diploma in Artificial
Intelligence; neither programme sponsored or supervised this study, which is reported here in
a personal capacity. Correspondence concerning this article should be addressed to Maximiliano
Speranza. E-mail: maximiliano.speranza@gmail.com.
The stopping criterion analyzed here was pre-registered and frozen with a DOI \emph{before}
any data were collected, which rules out post hoc selection of the criterion as an explanation
for the defect reported. The author declares no competing interests and received no funding
for this work. Code, stored validation histories, and the analysis script that reproduces every
number reported here are publicly available (see Open Practices).}

\abstract{%
Architecture comparisons in efficient sequence modeling commonly stop training with a rule of
the form ``improvement in validation accuracy over the last $w$ steps is below a tolerance
$\tau$,'' and then grant every architecture the resulting budget. We analyze one such rule that
was pre-registered and frozen with a DOI before data collection, and show that it could not
detect convergence---not because its tolerance was miscalibrated, but because the stopping
decision was taken below the signal-to-noise floor of the measurement. In a multi-query
associative recall benchmark, the delta-rule linear-attention condition improved at a rate of
$0.31$ accuracy points per evaluation in the regime treated as a plateau (positive in seven of
eight seeds; $t = 4.20$, $p = .004$), while the noise of the tested difference was $1.16$ points,
giving a signal-to-noise ratio of $0.27$. Under this ratio no tolerance separates the two errors
of interest: rules that detect a plateau also stop during genuine improvement. We show that a
proposed remedy---window means with a noise-calibrated tolerance---has a detection rate that is
algebraically fixed at $\Phi(2)$ regardless of window length, and would have declared convergence
while the model was still improving with probability $.73$. Decomposing the achievable gain, a
factor of $7.0$ comes from lengthening the comparison window and only $1.31$ from replacing the
two-point difference with a trend test. The failure is architecture-dependent: saturating
conditions exhibit noise three orders of magnitude smaller and close immediately, so the shared
budget covaries with the quantity under comparison. We give a power-based rule for choosing $w$
that can be pre-registered without pilot data. All analyses reuse stored validation histories
and require no additional training.}

\keywords{stopping criterion, early stopping, statistical power, pre-registration,
linear attention, associative recall, training budget, reproducibility}

\begin{document}
\maketitle

\section{Introduction}

Architecture comparisons in efficient sequence modeling rest on a decision that is rarely
examined: when to stop training. A common practice is to train a reference architecture until a
convergence rule fires and then grant every competing architecture the same number of steps
\parencite{li2019budgeted}. The procedure is sound only if convergence can be detected reliably.
If detection is unreliable, the shared budget inherits that unreliability, and so does the
comparison.

This article reports a case in which detection was impossible by construction, and characterizes
why. The setting is a pre-registered comparison of attention head compositions on multi-query
associative recall \parencite[MQAR;][]{arora2023zoology}. The stopping rule was specified,
frozen, and published with a DOI before any model was trained. That matters for the argument
developed here: the defect cannot be attributed to a rule chosen after seeing the results,
because the rule could not have been chosen after seeing them.

Our diagnosis differs from the one an inspection of the rule alone would suggest. The tolerance
was indeed smaller than the noise of the metric, and it is tempting to conclude that the rule
therefore measured noise instead of convergence. That conclusion does not survive contact with
the data. The regime the rule was applied to was not a plateau: the architecture of interest was
still improving, at a rate that a longer window would have detected. The rule failed because the
\emph{stopping decision} was taken below the signal-to-noise floor of the measurement, and no
choice of tolerance repairs a decision taken below that floor.

Our contribution is fourfold. First, we show that the regime treated as converged carried a
statistically detectable improvement trend, so the diagnosis must be framed in terms of
signal-to-noise rather than tolerance. Second, we show that recalibrating the tolerance to
measured noise---the natural remedy---produces a detection rate that is algebraically fixed and
independent of the window, and buys detection power at the cost of stopping during genuine
improvement. Third, we decompose the achievable gain and find that window length dominates the
choice of test statistic by a factor of five. Fourth, we show that the noise driving all of this
is architecture-dependent, which converts an apparently neutral procedural choice into a
confound aligned with the variable under study, and we give a power-based rule for $w$ that is
pre-registrable without pilot data.

\section{The Pre-Registered Setting}

The experiment compares three head compositions in a small transformer ($193{,}493$ parameters)
on MQAR: \texttt{softmax} (standard attention), \texttt{delta} (delta-rule linear attention),
and \texttt{mix22} (two softmax heads and two delta heads in the same layer). Delta-rule linear
attention updates a fixed-size associative state by correcting the current read before writing
\parencite{yang2024gated}; hybrid compositions are known to recover recall capacity that purely
linear variants do not reach \parencite{arora2023zoology,arora2024simple}.

Validation accuracy was evaluated every $500$ steps as the mean of top-1 accuracy over loads
$\{64, 96, 128\}$, on a validation set held fixed across evaluations by a fixed random seed. The
pre-registered rule, denoted D-004, declares convergence at step $N$ when
%
\begin{equation}
\mathrm{acc}(N) - \mathrm{acc}(N - w) < \tau ,
\label{eq:d004}
\end{equation}
%
with $w = 500$ and $\tau = 0.005$. Because evaluations are spaced $500$ steps apart,
Equation~\ref{eq:d004} compares two \emph{consecutive} evaluations. A condition closes its first
training phase only when all eight seeds satisfy Equation~\ref{eq:d004}; otherwise training
proceeds in blocks of $2{,}500$ steps up to a cap of $10{,}000$. The common budget for the
primary comparison is then the maximum of the per-condition convergence points, so one
condition's stopping behavior sets the budget granted to all.

\section{Method}

All analyses reuse validation histories stored during the original campaign; no model was
retrained. The data comprise eight \texttt{delta} runs, eight \texttt{mix22} runs, and three
\texttt{softmax} runs, plus one \texttt{delta} seed extended to $10{,}000$ steps and used only
where noted.

\subsection{Noise and Trend Estimation}

For each run we isolate the late-training regime---from step $4{,}000$ for \texttt{delta} and
$1{,}000$ for the saturating conditions---and fit an ordinary least-squares line to validation
accuracy against evaluation index. We estimate noise by two methods: the standard deviation of
residuals from that fit, and the standard deviation of successive differences divided by
$\sqrt{2}$. The second is unbiased under a linear trend and therefore serves as a check on the
first. The quantity that governs Equation~\ref{eq:d004} is the standard deviation of the tested
difference, $\sigma_D = \sigma\sqrt{2}$.

The trend is tested per run by the OLS slope $b$ against $H_0\colon b = 0$, one-tailed, and
across runs by a one-sample $t$ test on the eight slopes.

\subsection{Detection Probability}

We report two estimates. The parametric estimate models late-training observations as
$\mathrm{acc}(N) = \mu + \varepsilon$ with $\varepsilon \sim \mathcal{N}(0, \sigma^2)$, giving
$P = \Phi(\tau / \sigma_D)$ for one seed and $P^{k}$ for $k$ independent seeds. The
non-parametric estimate exploits the fact that Equation~\ref{eq:d004} with $w$ equal to the
evaluation interval \emph{is} a successive difference: we therefore report the observed fraction
of successive differences below $\tau$, which requires no distributional assumption and no
assumption of independence.

\section{Results}

\subsection{Metric Noise Is Architecture-Dependent}

Table~\ref{tab:ruido} reports late-training noise by condition. The two estimators agree to
within $1\%$, which also indicates that residual autocorrelation is negligible
(see Assumption Checks). The decisive quantity is $\tau / \sigma_D$: below $1$, the tolerance
cannot resolve the noise of the quantity it is compared against.

\begin{table}[htbp]
\caption{Late-Training Noise of Validation Accuracy by Condition}
\label{tab:ruido}
\small
\begin{tabular}{lccccc}
\toprule
Cond. & $n$ & $\sigma_{\text{resid}}$ & $\sigma_{\text{diff}}$ & $\sigma_D$ & $\tau/\sigma_D$ \\
\midrule
\texttt{delta}   & 8 & 0.00821 & 0.00829 & 0.01162 & 0.43 \\
\texttt{mix22}   & 8 & 0.00004 & 0.00004 & 0.00006 & 86.9 \\
\texttt{softmax} & 3 & 0.00004 & 0.00004 & 0.00006 & 88.8 \\
\bottomrule
\end{tabular}

\vspace{5pt}
{\small\raggedright \textit{Note.} $\sigma_D = \sigma_{\text{resid}}\sqrt{2}$ is the standard deviation of
the difference tested by Equation~\ref{eq:d004}. Estimates use eight late-training evaluations
per run for \texttt{delta} and four for the saturating conditions; the ratios for the latter
should be read as ``two orders of magnitude or more,'' not as precise values.
\par}
\end{table}

The three conditions differ in $\sigma$ by roughly three orders of magnitude. A single fixed
tolerance is thus highly informative for two of them and unusable for the third.

\subsection{The Regime Is Not a Plateau}

Before asking whether the rule detects convergence, we ask whether there was convergence to
detect. Table~\ref{tab:tend} tests the trend in the regime the rule was applied to.

\begin{table}[htbp]
\caption{Trend Test in the Regime Treated as Converged (Steps 4,000--7,500)}
\label{tab:tend}
\small
\begin{tabular}{lrrrrc}
\toprule
Seed & $b$ & $t$ & $p$ & Gain & D-004 \\
\midrule
0 & $+0.00259$ & 1.86 & .056 & $+1.81$ & yes \\
1 & $+0.00234$ & 1.72 & .068 & $+1.64$ & no \\
2 & $+0.00421$ & 2.54 & .022 & $+2.95$ & yes \\
3 & $+0.00334$ & 5.12 & .001 & $+2.34$ & no \\
4 & $-0.00014$ & $-0.12$ & .547 & $-0.10$ & yes \\
5 & $+0.00445$ & 3.58 & .006 & $+3.11$ & no \\
6 & $+0.00141$ & 1.37 & .110 & $+0.98$ & yes \\
7 & $+0.00680$ & 4.13 & .003 & $+4.76$ & no \\
\bottomrule
\end{tabular}

\vspace{5pt}
{\small\raggedright \textit{Note.} $b$ is the OLS slope per evaluation; $p$ values are one-tailed against
$H_0\colon b = 0$. ``Gain'' is the fitted improvement in accuracy points over the $3{,}500$ steps
of the window. The last column reports whether Equation~\ref{eq:d004} declared convergence at
step $7{,}500$.
\par}
\end{table}

The slope is positive in seven of eight seeds, with mean $+0.00312$ per evaluation; a one-sample
$t$ test on the eight slopes rejects zero, $t(7) = 4.20$, $p = .004$. Extrapolating the mean
slope predicts $+1.56$ accuracy points from step $7{,}500$ to $10{,}000$. The one seed actually
extended to $10{,}000$ gained $+1.39$ points of capacity at load~96 ($0.9167 \rightarrow
0.9306$), consistent with the projection.

This reframes the problem. The \texttt{delta} condition had not converged at $7{,}500$ steps.
Driving it to the step cap was not an artifact of a noisy rule; it was the correct outcome
reached for the wrong reason. The rule cannot be faulted for failing to detect a convergence
that had not occurred---but neither can it be credited, since Table~\ref{tab:tend} shows its
classifications bear no reliable relation to the measured trend.

\subsection{The Decision Is Taken Below the Signal-to-Noise Floor}

The two quantities can now be compared directly. The signal available to
Equation~\ref{eq:d004} is one evaluation interval of improvement, $b = 0.00312$; the noise of
the tested difference is $\sigma_D = 0.01162$. Their ratio is
%
\begin{equation}
\mathrm{SNR} = \frac{b}{\sigma_D} = 0.27 .
\label{eq:snr}
\end{equation}
%
This single number governs everything that follows. A stopping rule must balance two errors:
continuing to train after convergence, and stopping while improvement remains. With
$\mathrm{SNR} < 1$ the two error rates cannot be separated by any choice of $\tau$, because the
distributions being discriminated overlap almost completely. Tolerance selection is not the
lever; it merely slides a decision boundary along an axis on which the two hypotheses are not
distinguishable.

\subsection{Detection Probability}

For \texttt{delta}, the parametric model gives $P = \Phi(0.43) = .667$ per seed and
$P^{8} = .039$ for all eight simultaneously. The non-parametric estimate---the observed fraction
of successive differences below $\tau$ in the same regime, $N = 56$---gives $P = .554$,
95\% CI $[.423, .684]$, and $P^{8} = .009$. The non-parametric estimate is preferable: it
assumes nothing, and it fits the observation better. Four of eight seeds were classified as
converged; the expected count is $4.4$ under the non-parametric estimate ($P(X \leq 4) = .52$)
versus $5.3$ under the parametric one ($P(X \leq 4) = .26$).

Either way, the condition had less than a $4\%$ chance of ever closing its phase by
convergence---which is the correct outcome here, since it had not converged, but would have been
equally true had it converged.

\subsection{Recalibrating the Tolerance Does Not Help}

The natural remedy is to compare window means rather than points, and to set the tolerance from
measured noise:
%
\begin{equation}
\overline{\mathrm{acc}}_{\text{last } m} - \overline{\mathrm{acc}}_{\text{prev. } m} < c\,\sigma_D^{(m)},
\qquad \sigma_D^{(m)} = \sigma\sqrt{2/m} .
\label{eq:fix}
\end{equation}
%
With $m = 3$ and $c = 2$, the all-seed detection probability rises from $.039$ to $.832$. That
figure is not evidence of a better rule. Setting the tolerance to $c\,\sigma_D^{(m)}$ makes the
per-seed detection rate equal to $\Phi(c)$ \emph{identically}, for any $m$ and any $\sigma$: with
$c = 2$ it is $\Phi(2) = .977$ whatever the window. Table~\ref{tab:desc} separates the two
ingredients on our data. The entire gain comes from relaxing the tolerance, not from windowing.

\begin{table*}[htbp]
\caption{Decomposition of the Detection Gain}
\label{tab:desc}
\begin{tabular}{lccc}
\toprule
Variant & $\tau$ & $\sigma_D$ & $P^{8}$ \\
\midrule
D-004 as pre-registered              & 0.0050 & 0.01162 & .039 \\
Windowing only ($m = 3$, $\tau$ fixed) & 0.0050 & 0.00671 & .126 \\
Retuned tolerance only ($2\sigma_D$) & 0.0232 & 0.01162 & .832 \\
Equation~\ref{eq:fix} in full        & 0.0134 & 0.00671 & .832 \\
\bottomrule
\end{tabular}

\vspace{5pt}
{\small\raggedright \textit{Note.} $P^{8}$ is the probability that all eight seeds are simultaneously
declared converged while on a plateau. The two lower rows coincide because setting
$\tau = 2\sigma_D$ fixes the per-seed rate at $\Phi(2)$ irrespective of the window.
\par}
\end{table*}

The cost of that relaxation is the error the detection rate does not measure. Under a true
improvement rate $b$, the probability that Equation~\ref{eq:fix} declares convergence anyway is
$\Phi\!\left(c - b\,m/\sigma_D^{(m)}\right)$. At the improvement rate actually measured, that
probability is $.73$: the proposed rule would have stopped \texttt{delta} while it was still
gaining roughly one accuracy point per window. Its minimum detectable improvement, at a $5\%$
false-stop rate, is $4.07$ points per $2{,}500$-step block---larger than the $1.39$-point gain
that the extended seed actually realized over such a block. A rule tuned to detect plateaus
reliably in this regime is, necessarily, a rule that stops during real progress.

\subsection{Window Length Dominates the Test Statistic}

What does raise the floor is the window. For $n$ equally spaced evaluations, the OLS slope has
standard error
%
\begin{equation}
\mathrm{SE}(b) = \sigma \sqrt{\frac{12}{n(n^{2}-1)}} ,
\label{eq:se}
\end{equation}
%
which falls as $n^{-3/2}$, against $n^{-1/2}$ for averaging. On our data, with $n = 8$
evaluations, $\mathrm{SE}(b) = 0.00127$ and the measured mean slope yields $t = 2.46$: the same
eight points that leave Equation~\ref{eq:d004} at $\mathrm{SNR} = 0.27$ support detection when
used jointly. The total gain of $9.2\times$ decomposes into $7.0\times$ from comparing points
seven evaluations apart instead of one, and $1.31\times$ from using the trend statistic instead
of the endpoint difference. Window length, not the choice of statistic and not the tolerance, is
the dominant term.

This yields a rule that can be fixed in advance. Given a smallest improvement rate $b^{*}$ that
the study commits to not overlooking, a one-tailed level $\alpha$, and power $1 - \beta$, the
window must satisfy
%
\begin{equation}
b^{*} \geq (z_{\alpha} + z_{\beta})\, \sigma \sqrt{\frac{12}{n(n^{2}-1)}} .
\label{eq:power}
\end{equation}
%
Table~\ref{tab:pot} evaluates Equation~\ref{eq:power} at the noise level measured here. It shows
what the original design would have required: detecting an improvement of one accuracy point per
$2{,}500$ steps needs about eleven evaluations, whereas D-004 used two.

\begin{table}[htbp]
\caption{Detectable Improvement as a Function of Window Length}
\label{tab:pot}
\small
\begin{tabular}{cccc}
\toprule
$n$ & Window & $\mathrm{SE}(b)$ & $b^{*}$ \\
\midrule
4  & 1,500 & 0.00367 & 4.56 \\
6  & 2,500 & 0.00196 & 2.44 \\
8  & 3,500 & 0.00127 & 1.58 \\
10 & 4,500 & 0.00090 & 1.12 \\
12 & 5,500 & 0.00069 & 0.85 \\
16 & 7,500 & 0.00045 & 0.55 \\
20 & 9,500 & 0.00032 & 0.40 \\
\bottomrule
\end{tabular}

\vspace{5pt}
{\small\raggedright \textit{Note.} $n$ = number of evaluations in the window; ``Window'' is its span in
training steps. $\sigma = 0.00821$, $\alpha = .05$ one-tailed, power $= .80$. $b^{*}$ is the
smallest improvement rate detectable with that power, in accuracy points per $2{,}500$ steps.
\par}
\end{table}

Equation~\ref{eq:power} needs $\sigma$, which is not known before training. Two routes avoid
circularity. Where the validation set is resampled at each evaluation, sampling error gives an
analytic lower bound, $\sigma \geq \sqrt{p(1-p)/n_{\text{val}}}$, computable at design time.
Where it is held fixed, as here, that term vanishes and the residual noise is optimizer
variability, which must be measured; a short pilot on any one condition suffices, and---unlike
retuning a tolerance after seeing results---it fixes a window, not a verdict.

\subsection{Propagation to the Training Budget}

Because the rule sets the common budget, its architecture-dependence is not neutral. Saturating
conditions have $\tau/\sigma_D \approx 87$ and close at the first opportunity; the non-saturating
condition is still improving and is driven to the cap. Whichever budget is then adopted is
either insufficient for the condition that is still learning or far beyond what the saturating
ones need, and \emph{which of the two occurs is decided by the conditions that saturate}. The
granted budget thus covaries with task performance---the very quantity under comparison.

The magnitude is material. The pre-registered decision margin for the primary hypothesis was
$2$ accuracy points; the seed extended from $7{,}500$ to $10{,}000$ steps gained $1.39$ points of
capacity at load~96. A budget-induced difference of the same order as the decision margin is not
a rounding concern. This estimate rests on a single extended seed and is indicative rather than
conclusive; the trend analysis of Table~\ref{tab:tend}, which uses all eight, supports the same
direction and magnitude.

\subsection{Assumption Checks}
\label{sec:assump}

\paragraph{Residual autocorrelation.} Lag-1 autocorrelation of detrended residuals is negative
in seven of eight runs, median $-0.315$, which invites a correction to $\mathrm{Var}(D)$. The
correction would be spurious. With $n = 8$ points and a fitted trend removed, the lag-1 estimator
is strongly biased downward even for white noise: simulation under white noise ($20{,}000$
replicates) gives $E[\hat\rho_1] = -0.287$, 95\% interval $[-0.803, +0.310]$, placing the
observed median at the $49$th percentile of the null. This is the small-sample bias of
autocorrelation estimators \parencite{marriott1954bias}. The data contain an independent check:
were $\rho_1 = -0.315$ real, $\sigma_{\text{diff}}/\sigma_{\text{resid}}$ would equal
$\sqrt{1-\rho_1} = 1.147$; the observed ratio is $1.010$, implying $\rho_1 = -0.02$. We therefore
treat the residuals as serially uncorrelated. We note the point because the tempting
correction runs in the direction that flatters the argument, and would have been reported as
support for it.

\paragraph{Source of the noise.} The validation set is held fixed across evaluations by a fixed
seed, and comprises roughly $37{,}000$ queries; resampling error would contribute $\sigma \approx
0.0018$ and, being fixed, contributes none. The measured $\sigma = 0.0082$ is therefore
attributable to optimizer variability, not to measurement error. Enlarging the validation set
would not reduce it; temporal averaging or weight averaging would.

\paragraph{Choice of regime onset.} Varying the onset from step $2{,}500$ to $5{,}000$ leaves
$\tau/\sigma_D$ between $0.39$ and $0.51$, always well below unity, and leaves the sign and
significance of the trend unchanged.

\paragraph{Does validation accuracy track the quantity of interest?} Across the eight
\texttt{delta} seeds, final validation accuracy correlates with capacity at load~96 at
$r = .96$ and at load~128 at $r = .99$. The metric is not blind to the axis under study; it is
noisy relative to its own rate of change. The problem is one of power, not of construct.

\section{Discussion}

\subsection{Implications}

The practical cost of a stopping rule below the noise floor is not primarily wasted computation,
although in our campaign roughly $80\%$ of the projected budget was attributable to it. The
serious cost is that the training budget becomes a function of the metric's noise, which is
itself a function of where the architecture sits on the task. Architectures that saturate stop
early; architectures in the interesting intermediate regime are driven to the cap. Any
comparison drawn across that asymmetry is confounded, and the confound is aligned with the
variable the study exists to measure.

The episode also carries a lesson about pre-registration. Freezing a stopping rule is good
practice, but a rule frozen without a power analysis freezes an unexamined assumption about the
noise of the metric. Our pre-registration fixed $\tau$ and $w$ and said nothing about $\sigma$;
Equation~\ref{eq:power} is what should have been registered instead, since it states what the
study commits to detecting rather than how it will compute a difference. This is the analogue,
for training-time decisions, of registering a target effect size rather than a $p$-value
threshold.

\subsection{Relation to Prior Work}

That validation-based early stopping is unstable in noisy regimes is established
\parencite{prechelt1998automatic}, as is the need to control training budget when comparing
architectures \parencite{li2019budgeted}, and seed-sampling noise has been shown to affect the
reliability of synthetic benchmarks \parencite{chanin2026reliable}. The specific contribution here
is the chain that links them and its quantification: a stopping decision taken below the
signal-to-noise floor, a floor that differs by three orders of magnitude across the architectures
being compared, and a budget rule that transmits that difference into the comparison. We also
report, as a cautionary detail, that the two natural repairs---retuning the tolerance to measured
noise, and averaging over a window---address the smaller term while leaving the dominant one
untouched.

\subsection{Limitations}

This is a single documented case, not a general study. It involves one task family (MQAR), one
model scale ($193$K parameters), one hardware platform, three conditions, and eight seeds per
condition. We do not claim that fixed-tolerance rules fail generally, nor that any published
result is invalidated by this mechanism. The budget-effect estimate at load~96 rests on one
extended seed; the trend analysis that corroborates it uses all eight but extrapolates beyond
the observed range. The power rule of Equation~\ref{eq:power} assumes an approximately linear
improvement trend and homoscedastic noise over the window, which is an approximation for
late-training dynamics and would fail near a sharp transition. Finally, Equation~\ref{eq:power}
is derived here and validated only retrospectively; prospective validation remains future work,
and is the natural next step for this line.

\subsection{Recommendations}

We suggest four inexpensive practices. First, before fixing a stopping rule, state the smallest
improvement rate the study commits to detecting, and size the window from it via
Equation~\ref{eq:power}; report the resulting power. Second, report $\tau/\sigma_D$ rather than
$\tau$ alone, since the tolerance is uninterpretable without the noise of the quantity it is
compared against. Third, report noise per condition: a rule adequate for one architecture may be
unusable for another in the same study, and that asymmetry is a confound, not an inconvenience.
Fourth, when a rule is relaxed to raise its detection rate, report the complementary error---the
probability of stopping under the smallest improvement that matters---because detection rate
alone can always be driven to any desired value.

\section{Conclusion}

A stopping rule compares an improvement against a tolerance, but what determines whether it can
work is neither of those quantities in isolation: it is the ratio between the improvement
accumulated over the comparison window and the noise of the comparison. When that ratio is
below unity, the rule reports which way the last evaluation happened to fall, and no tolerance
recovers the decision. Lengthening the window does, at a rate that makes the fix cheap. When such
a rule also sets the shared training budget of an architecture comparison, its
architecture-dependent noise sensitivity converts a procedural choice into a confound aligned
with the variable under study. Diagnosing this required nothing beyond validation histories that
were already being stored, and correcting it requires no additional training.

\section*{Open Practices}

The pre-registration was frozen and published before data collection
(DOI:~10.5281/zenodo.21495252). Code, stored validation histories, and the analysis script that
reproduces the reported numbers are available at
\url{https://github.com/SperanzaMax/telar-ligamento}; the analysis runs on CPU in under a
minute. One exception must be stated: the single \texttt{delta} seed extended to $10{,}000$
steps, cited once for its capacity gain of $1.39$ points, was lost to a session failure before
it could be versioned, and is the only quantity here that cannot be recomputed from the
repository. Every other figure, including the trend analysis that independently predicts that
gain, derives from the eight versioned runs.

\printbibliography

\end{document}
